% ===================================================================
%  TABELL Nr. 5 — D'Haus a säi Bau.
%  Traduction 2026 d'apres texto/io, controlee sur texto/fr, texto/de
%  et texto/nl. Consigne : voir texto/lb/10-tabell-01.tex.
%
%  LE CHANTIER EST ALLEMAND ET LA MAISON EST FRANCAISE, et c'est la
%  demonstration la plus nette de cette colonne.
%
%  Les hommes qui batissent portent tous un nom germanique :
%  « Maurer » le macon, « Zimmermann » le charpentier, « Schräiner »
%  le menuisier, « Dachdecker » le couvreur, « Spengler » le
%  zingueur, « Schlässer » le serrurier, « Moler » le peintre,
%  « Glaser » le vitrier, « Steemetzer » le tailleur de pierre,
%  « Handlanger » le manoeuvre, « Gäertner » le jardinier.
%
%  Les pieces qu'ils ont faites portent presque toutes un nom
%  francais : « Vestibule », « Corridor », « Salon », « Bureau »,
%  « Office », « Veranda », « Palier », « Alkove », « Terrass »,
%  « Mansarden », « Parterre ». Ne restent germaniques que les trois
%  pieces ou l'on ne recoit personne : « Kichen » la cuisine,
%  « Schlofzëmmer » la chambre a coucher, « Buedzëmmer » la salle de
%  bains.
%
%  L'HOMME QUI TRAVAILLE EST NOMME DANS LA LANGUE DU TRAVAIL, LA
%  PIECE OU L'ON RECOIT DANS LA LANGUE DE LA SOCIETE. C'est la mesure
%  double du tableau 2 — deux voisins, deux distances qui varient en
%  sens contraire — appliquee cette fois a un SEUL BATIMENT, et l'on
%  y voit la ligne passer entre la cour et le salon.
%
%  ET LE LUXEMBOURGEOIS N'AJOUTE RIEN A LA NOTE DU FEUILLET 35, ce
%  qui ne lui etait pas arrive. La note enumere les mots nationaux
%  qui rendent le francais « solive » : le tcheque y avait mis
%  « stropnice », le lituanien « sija ». Le luxembourgeois dit
%  « Balken », c'est-a-dire exactement le mot allemand que la note
%  cite deja. On l'ecrit tout de meme, et c'est la seule ligne du
%  volume ou deux langues de la liste donnent le meme mot.
% ===================================================================

%%K t05-apar-1 apar
\VUsaut{6.0mm}
\VUtitre{15.0pt}{60}{TABELL Nr. 5}
\VUsaut{1.4mm}
\VUtitre{11.4pt}{0}{D'Haus a säi Bau.}
\VUsaut{2.2mm}

%%K t05-tit-1 sub
\VUpk{10.2pt}{40}{Gutt Begéinung.}
\VUsaut{3.0mm}

%%K t05-01-1 p
\textsc{Jang}. --- Monni, ech géif ganz gär wëssen, wéi een en \VUgras{Haus} baut.

%%K t05-01-2 p
\textsc{De Monni} (80). --- Dat ass ganz einfach, mäi Frënd, komm mat mir, ech weisen
der dat, wat den Här Bernard (79) baue léisst a wat ech bewunne wäert.

%%K t05-01-3 p
\textsc{Jang}. --- A wien huet dann de \VUgras{Plang} \textsuperscript{(76)} vun dësem
Haus gezeechent?

%%K t05-01-4 p
\textsc{De Monni}. --- Den \VUgras{Architekt} \textsuperscript{(75)}; hien huet eng
ganz kloer Zeechnung virgeluecht, déi ugeholl gouf, an den
\VUgras{Entrepreneur} \textsuperscript{(77)} huet d'Aarbechten ugefaangen.

%%K t05-01-5 p
\textsc{Jang}. --- A wien ass dann den Entrepreneur?

%%K t05-01-6 p
\textsc{De Monni}. --- Dat ass den Här Wäiss, de Papp vun dengem Frënd Jhemp. Hie
leet d'Aarbechter zesumme mam Här Rout, dem Architekt.

%%K t05-01-7 p
Kuck! Do kënnt grad rechtzäiteg den Här Wäiss mat sengem \VUgras{Lineal}
\textsuperscript{(78)}, an den Här Rout mat sengem Plang.

%%K t05-01-8 p
Gudde Moien, d'Hären. Wéi geet et Iech?

%%K t05-01-9 p
\textsc{Den Här Rout}. --- Ganz gutt, merci. Gudde Moien, Jang, wéi ass dëst Kand
gewuess!

%%K t05-01-10 p
\textsc{De Monni}. --- Jo, wierklech, hien ass e klenge Mann. Hien ass ganz eescht,
ee muss him alles erklären, wat hie gesäit. Haut géif hie gär wëssen, wéi een en Haus
baut.

%%K t05-tit-2 sub
\VUpk{10.2pt}{40}{D'Aarbechter.}
\VUsaut{3.0mm}

%%K t05-01-11 p
\textsc{Den Här Wäiss}. --- Also komm mat mir, klenge Frënd, ech erklären der dat
gläich, well ech hunn d'Kanner ganz gär, déi sech wëllen instruéieren.

%%K t05-01-12 p
Du gesäis deen Aarbechter, deen eng \VUgras{Schëpp} \textsuperscript{(82)} an der
Hand hält: dat ass en \VUgras{Äerdaarbechter} \textsuperscript{(81)}. Hie gruewt eng
\VUgras{Grouf} \textsuperscript{(46)}, fir d'Waasserleedung ze leeën. Hien huet och de
Buedem op der Plaz ausgegruewen, wou d'Haus steet, fir datt de Maurer déi véier
\VUgras{Maueren} opriichte kann. Deen Deel vun dëse Maueren, deen am Buedem ass,
heescht d'\VUgras{Fundamenter} \textsuperscript{(2)}. De Raum tëscht de Fundamenter
bilt de \VUgras{Keller} \textsuperscript{(3)}. Dëse Keller huet e klengt Fënsterchen,
dat \VUgras{Kellerfënster} \textsuperscript{(5)} heescht, eng \VUgras{Dier}
\textsuperscript{(4)} an eng Trap.

%%K t05-01-13 p
De \VUgras{Maurer} \textsuperscript{(108)} huet d'Maueren weidergebaut an
d'Ouverturë fir d'\VUgras{Dieren} \textsuperscript{(8)} an d'\VUgras{Fënsteren}
\textsuperscript{(17)} gelooss.

%%K t05-01-14 p
\textsc{Jang}. --- Mä dee Maurer gesinn ech net!

%%K t05-01-15 p
H. \textsc{Wäiss}. --- Elo huet hien am Haus fäerdeg geschafft. Hien ass beim
\VUgras{Päerdsstall} \textsuperscript{(54)}, op sengem \VUgras{Gerüst}
\textsuperscript{(110)}, an hie baut d'Mauer vun der \VUgras{Wäscherei}
\textsuperscript{(56)}.

%%K t05-01-16 p
Hien huet eng Kell, eng Mierdelkëscht an e \VUgras{Bläilot} \textsuperscript{(109)}.
Mä du wäerts der dës Nimm net behalen.

%%K t05-01-17 p
\textsc{Jang}. --- Dach jo, well ech froen elo mäi Monni ëm eng kleng Kell an eng
Mierdelkëscht, an e Bläilot maachen ech mer selwer mat engem Fuedem.

%%K t05-tit-3 sub
\VUsaut{3.0mm}
\VUpk{10.2pt}{40}{D'Material.}
\VUsaut{3.0mm}

%%K t05-01-18 p
\textsc{De Monni}. --- Ah! ganz gutt. Mä sos mer, Jang, wat brauch ee fir en Haus ze
bauen?

%%K t05-01-19 p
\textsc{Jang}. --- Ee brauch \VUgras{Quadersteng} \textsuperscript{(59)} a
\VUgras{Mierdel} \textsuperscript{(65)}.

%%K t05-01-20 p
\textsc{De Monni}. --- Richteg, kuck dee Mann do, mat dem groussen Hutt, op sengem
\VUgras{Kärchen} \textsuperscript{(136)}. Dat ass e \VUgras{Fuhrmann}
\textsuperscript{(135)}. All Dag bréngt hien hei \VUgras{Steeblöck}
\textsuperscript{(60)}. Hie lued säi Won am Haff of, an d'\VUgras{Steemetzer}
\textsuperscript{(83)} haen d'Blöck a segen se mat hirer \VUgras{Steesäg}
\textsuperscript{(85)}. En \VUgras{Handlanger} \textsuperscript{(146)} dréit se dem
Maurer, deen se setzt.

%%K t05-01-21 p
Mëttlerweil bereet de \VUgras{Mierdelmaacher} \textsuperscript{(87)} de Mierdel vir.
Hie brauch \VUgras{Sand} \textsuperscript{(62)}, \VUgras{Kallek} \textsuperscript{(63)}
a \VUgras{Waasser} \textsuperscript{(64)}. De Kallek ass nieft him an engem
\VUgras{Sak} \textsuperscript{(71)}, e \VUgras{Maurergesell} \textsuperscript{(111)}
bréngt him de Sand op enger \VUgras{Schubkar} \textsuperscript{(112)}, an de
\VUgras{Léierbouf} \textsuperscript{(113)} ass bei der Pomp (58). Hie fëllt en
\VUgras{Eemer} \textsuperscript{(114)}. De Mierdel ass geschwënn prett. De
\VUgras{Mierdeldréier} \textsuperscript{(107)} hëlt en a dréit en iwwer d'Leeder erop
op d'Gerüst, wou e Maurer schafft, deen déi Säit vum Haus fäerdeg mécht.

%%K t05-01-22 p
An elo, wéi heescht den Aarbechter, deen um \VUgras{Daach} \textsuperscript{(31)}
schafft?

%%K t05-01-23 p
\textsc{Jang}. --- Dat ass en \VUgras{Dachdecker} \textsuperscript{(118)}.

%%K t05-tit-4 sub
\VUsaut{3.0mm}
\VUpk{10.2pt}{40}{D'Daach vum Haus.}
\VUsaut{3.0mm}

%%K t05-01-24 p
H. \textsc{Wäiss}. --- Jo, mä fir d'éischt kënnt den \VUgras{Zimmermann}
\textsuperscript{(115)}.

%%K t05-01-25 p
Den Zimmermann schafft de \VUgras{Dachstull} \textsuperscript{(35)}. Hie setzt
d'\VUgras{Balken} \textsuperscript{(37)} mat \VUgras{Zapfen} \textsuperscript{(117)}
zesummen. Déi Aarbechter do uewen, mat hire breede Samtäermel, sinn Zëmmerleit. Deen
een hält e Bäckelchen, an deen anere schneit en Zapfen mat senger \VUgras{Aaxt}
\textsuperscript{(116)}. Deen, deen nieft dem \VUgras{Kamäin} \textsuperscript{(41)}
ass, ass den Dachdecker. Hien neelt Latten op déi (*) Balken, an \VUgras{Schifer}
\textsuperscript{(66)} op d'Latten. Heiansdo leet hien Zillen.

%%K t05-noto-1 noto
\VUnotes{143.2mm}{%
(*) \VUgras{Balk-o} = däitsch \textit{Balken}, engl. \textit{joist}, franz.
\textit{solive}, ital. \textit{travicello}, russ. \textit{balka}, lëtz.
\textit{Balken}, span. a port. \textit{viga}. Eng technesch Wurzel, déi bis elo net
ugeholl ass, mä déi virzeschloen ass, fir de Sënn vum franséische \textit{solive}
erëmzeginn, dat weder e \textit{trabo} (franz. \textit{poutre}) nach e Bäckelchen
ass. Et ass e viereckeg behaue Holzstaang, deen de Buedem an d'Plafong dréit.%
}

%%K t05-01-26 p
D'Daach vum Päerdsstall ass \VUgras{mat Zille gedeckt} \textsuperscript{(55)}, dat
vum Haus ass \VUgras{mat Schiefer gedeckt} \textsuperscript{(32)}, an dat vun der
Veranda ass aus \VUgras{Zink} \textsuperscript{(24)}.

%%K t05-01-27 p
\textsc{Jang}. --- Schafft den Dachdecker och d'Zinkdiecher?

%%K t05-01-28 p
H. \textsc{Wäiss}. --- Nee, mä de \VUgras{Spengler} \textsuperscript{(121)} oder de
\VUgras{Bläispengler,} deen d'\VUgras{Dachfënster} \textsuperscript{(38)} op d'Daach
vum Haus setzt. Hie setzt och d'\VUgras{Reenoflafrouer} \textsuperscript{(43)} an
d'\VUgras{Dachrënnen} \textsuperscript{(42)}. Hie léit d'Zink an d'Blech zesummen.
Hien huet de \VUgras{Blëtzofleeder} \textsuperscript{(40)} an d'\VUgras{Wiederhunn}
\textsuperscript{(39)} um \VUgras{Giewel} \textsuperscript{(36)} befestegt.

%%K t05-01-29 p
\textsc{Jang}. --- Ass d'Haus dann esou fäerdeg?

%%K t05-tit-5 sub
\VUsaut{3.0mm}
\VUpk{10.2pt}{40}{D'Handwierksgeschir.}
\VUsaut{3.0mm}

%%K t05-01-30 p
H. \textsc{Wäiss}. --- Net ganz. De \VUgras{Gipser} \textsuperscript{(99)} baut mat
\VUgras{Zillen} \textsuperscript{(61)} d'Trennmaueren, déi et a verschidde Stuffen
deelen. Hie verputzt mat \VUgras{Gips} \textsuperscript{(70)} d'\VUgras{Plafongen}
\textsuperscript{(26)} an d'Maueren. Elo schafft hien d'Plafong vun der
\VUgras{Veranda} \textsuperscript{(33)}. Hien huet, wéi de Maurer, eng \VUgras{Kell}
\textsuperscript{(100)} an eng \VUgras{Mierdelkëscht} \textsuperscript{(101)}.

%%K t05-01-31 p
Duerno schafft de Schräiner d'Dieren, d'Fënsteren, de \VUgras{Parquet}
\textsuperscript{(16)} an d'\VUgras{Fënsterlueden} \textsuperscript{(19)}.

%%K t05-01-32 p
Elo schafft hien am Haff op senger \VUgras{Wierkbänk} \textsuperscript{(90)}. Hien
huet eng \VUgras{Säg} \textsuperscript{(94)}, fir d'Holz (69) ze segen, en
\VUgras{Hiwwel} \textsuperscript{(91)}, e \VUgras{Schrauwstack}
\textsuperscript{(92)}, e \VUgras{Beitel} \textsuperscript{(94 bis)}, en
\VUgras{Zirkel} \textsuperscript{(95)} an en hëlzenen \VUgras{Hummer}
\textsuperscript{(95 bis)}.

%%K t05-01-33 p
\textsc{Jang}. --- Ah! mä et ass méi ameséiereg, ze schräineren wéi ze mauren. Monni,
kaafs du mer eng Wierkbänk, eng Säg an en Hiwwel? Well ech maache geschwënn eng Hütt
fir de Medor.

%%K t05-01-34 p
\textsc{De Monni}. --- Jo, klenge Bouf.

%%K t05-01-35 p
\textsc{Jang}. --- Wéi zefridde sinn ech! A wien ass deen, dee bei der Entréesdier
steet?

%%K t05-tit-6 sub
\VUsaut{3.0mm}
\VUpk{10.2pt}{40}{D'Molen.}
\VUsaut{3.0mm}

%%K t05-01-36 p
H. \textsc{Wäiss}. --- Dat ass e \VUgras{Moler} \textsuperscript{(102)}. Hie steet op
senger Leeder mat engem Dëppe \VUgras{Faarf} \textsuperscript{(103)} a mat sengem
\VUgras{Pinsel} \textsuperscript{(104)}. Hie molt d'Holz vun der Entréesdier, fir datt
et méi laang hält.

%%K t05-01-37 p
Hie pecht och d'Tapéiten an de Wunnengen.

%%K t05-01-38 p
Den \VUgras{Tapisséier} \textsuperscript{(107)}, deen op enger \VUgras{Leeder}
\textsuperscript{(106)} um \VUgras{Balkon} \textsuperscript{(28)} steet, setzt e
\VUgras{Store} \textsuperscript{(29)}.

%%K t05-01-39 p
Den Aarbechter, deen nieft der grousser Fënster steet, ass de \VUgras{Glaser}
\textsuperscript{(96)}. Hie befestegt d'\VUgras{Fënsterschäiwen}
\textsuperscript{(18)} mat \VUgras{Kitt} \textsuperscript{(73)}. Deen aneren
Aarbechter, dee bei der Kellerdier knéit, ass e \VUgras{Schlässer}
\textsuperscript{(97)}. Hie setzt e \VUgras{Schlass} \textsuperscript{(6)}. Seng
\VUgras{Feil} \textsuperscript{(98)} läit nieft him.

%%K t05-01-40 p
\textsc{Jang}. --- Ass deen Aarbechter, dee mat sengem \VUgras{Hummer}
\textsuperscript{(84)} op Eisen kläppt, och e Schlässer?

%%K t05-01-41 p
H. \textsc{Wäiss}. --- Méi genee gesot, dat ass e Schmadd (125). Hie schmitt
\VUgras{Eisen} \textsuperscript{(67)} fir den \VUgras{Elektriker}
\textsuperscript{(124)}, dee um Daach vun der \VUgras{Remise} \textsuperscript{(51)}
ass. Hien huet e \VUgras{Schmiddefeier} \textsuperscript{(126)}, en \VUgras{Amboss}
\textsuperscript{(127)} an eng \VUgras{Zaang} \textsuperscript{(128)}.

%%K t05-01-42 p
Den Elektriker befestegt de \VUgras{koperen} \textsuperscript{(44)} \VUgras{Drot} vum
Telefon.

%%K t05-tit-7 sub
\VUpk{10.2pt}{40}{D'Gaardenaarbecht.}
\VUsaut{3.0mm}

%%K t05-01-43 p
\textsc{De Monni}. --- Hannen am \VUgras{Haff} \textsuperscript{(45)}, nieft dem
\VUgras{Gitter} \textsuperscript{(47)} an dem \VUgras{Portal} \textsuperscript{(48)},
gesäis du de \VUgras{Gäertner} \textsuperscript{(129)}. Hie bereet de
\VUgras{Gäertchen} \textsuperscript{(49)} vir. Hien huet de Buedem mat senger
\VUgras{Spuet} \textsuperscript{(131)} ëmgestach, hie séit d'Somen a rechent mat dem
\VUgras{Rechen} \textsuperscript{(130)}. Duerno wäert hie mat senger
\VUgras{Waasserkan} \textsuperscript{(132)} d'\VUgras{Bieder} \textsuperscript{(150)}
begéissen, an d'Planze wuessen.

%%K t05-01-44 p
De \VUgras{Päerdsknecht} \textsuperscript{(133)}, deen d'Päerd grad gepflegt an him
Fudder ginn huet, botzt d'\VUgras{Kutsch} \textsuperscript{(52)} mat engem Schwamm,
enger \VUgras{Handpomp} \textsuperscript{(134)} an engem Eemer Waasser. Duerno féiert
hien se nees an d'Remise a mécht d'Dier mat dem décke \VUgras{Rigel}
\textsuperscript{(53)} zou.

%%K t05-01-45 p
Elo bass du zefridden, mengen ech, du hues genuch Erklärunge kritt.

%%K t05-01-46 p
\textsc{Jang}. --- Jo, Monni, mä ech géif gär d'Bannescht vum Haus gesinn.

%%K t05-01-47 p
\textsc{De Monni}. --- Dat gëtt muer. Den Här Rout wäert der de Plang erkläre an der
den Numm vun all Stuff soen.

%%K t05-tit-8 sub
\VUsaut{3.0mm}
\VUpk{10.2pt}{40}{Déi bannescht Antelung vum Haus.}
\VUsaut{2.4mm}

%%K t05-apar-2 apar
{\centering\textit{(Kuck de Plang.)}\par}
\VUsaut{2.4mm}

%%K t05-01-48 p
\textsc{Den Architekt}. --- Komm, Jang, ech loossen dech elo déi verschidden Deeler
vum Haus entdecken.

%%K t05-01-49 p
Gi mer an d'\VUgras{Parterre} \textsuperscript{(7)} eran. Do ass de Knäppchen vun
der \VUgras{Schell} \textsuperscript{(11)}. Mir klamme mat dräi \VUgras{Stufen}
\textsuperscript{(14)} op de \VUgras{Perron} \textsuperscript{(9)}, mir iwwerschreiden
d'\VUgras{Schwell} \textsuperscript{(10)} vun der Entrée-\VUgras{Dier}
\textsuperscript{(8)}, a do stinn mer um Fouss vun der \VUgras{Trap}
\textsuperscript{(13)}.

%%K t05-01-50 p
\textsc{Jang}. --- Dat ass de Corridor.

%%K t05-01-51 p
\textsc{Den Architekt}. --- Nee, dat ass de \VUgras{Vestibule}
\textsuperscript{(12)}. De \VUgras{Corridor} \textsuperscript{(a)} ass hannen. Riets,
do sinn de \VUgras{Salon} \textsuperscript{(b)} an d'\VUgras{Iesszëmmer}
\textsuperscript{(c)}, dem Iesszëmmer géintiwwer sinn d'\VUgras{Kichen}
\textsuperscript{(d)} an d'\VUgras{Office} \textsuperscript{(e)}, an duerno vir de
\VUgras{Bureau} \textsuperscript{(f)}. D'\VUgras{Veranda} \textsuperscript{(23)} mat
hirer \VUgras{Glasdier} \textsuperscript{(25)} ass nieft dem Salon.

%%K t05-01-52 p
Elo klamme mer d'Trap erop. Halt dech un dem \VUgras{Handleef}
\textsuperscript{(15)} a ziel d'Stufen. Et sinn der véierzéng. Do ass de
\VUgras{Palier} \textsuperscript{(m)}, mir sinn um éischte \VUgras{Stack}
\textsuperscript{(27)}.

%%K t05-tit-9 sub
\VUsaut{3.0mm}
\VUpk{10.2pt}{40}{Den éischte Stack.}
\VUsaut{3.0mm}

%%K t05-01-53 p
Der Trap géintiwwer sinn d'\VUgras{Toilette} \textsuperscript{(20)} an
d'\VUgras{Kabinett} \textsuperscript{(j)}. Riets e schéint \VUgras{Schlofzëmmer}
\textsuperscript{(g)} mat zwou Fënsteren op d'\VUgras{Fassad} \textsuperscript{(1)}
vum Haus. Lénks sinn d'\VUgras{Buedzëmmer} \textsuperscript{(1)} an
d'\VUgras{Frëndenzëmmer} \textsuperscript{(i)} mat enger grousser \VUgras{Alkove}
\textsuperscript{(h)}: nieft dem Schlofzëmmer ass d'\VUgras{Kannerzëmmer}
\textsuperscript{(k)}. Et geet op eng breet \VUgras{Terrass} \textsuperscript{(21)}.
Ënner dëser Terrass ass eng aner Toilette (20), an un der Mauer, do ass
d'\VUgras{Klack} \textsuperscript{(22)}, déi fir d'Iesse laut.

%%K t05-01-54 p
\textsc{Jang}. --- Gi mer op den \VUgras{zweete Stack}.

%%K t05-01-55 p
\textsc{Den Architekt}. --- Et gëtt kee zweete Stack. Ënnert dem Daach sinn nëmme
\VUgras{Mansarden} \textsuperscript{(33)} an de Späicher (34). Mir hunn d'Visite
fäerdeg, mir kënne nees erofklammen.

%%K t05-01-56 p
\textsc{Jang}. --- Merci, Här, dat ass ganz interessant. Ech erzielen elo mengem Papp
alles, wat ech gesinn hunn.

\VUsaut{4.0mm}
\VUfilet{22mm}
